
\documentclass[11pt]{article}

% ============================================================
% Cover & Main Document Setup
% ============================================================
\usepackage[utf8]{inputenc}
\usepackage[T1]{fontenc}
\usepackage[margin=1in]{geometry}
\usepackage{graphicx}
\usepackage{tabularx}
\usepackage{booktabs}
\usepackage{array}
\usepackage{amsmath,amssymb}
\usepackage[round]{natbib}
\usepackage{enumitem}
\usepackage{xcolor}
\usepackage{hyperref}
\hypersetup{
    colorlinks=true,
    linkcolor=blue!50!black,
    citecolor=blue!50!black,
    urlcolor=blue!50!black
}

\newcommand{\todo}[1]{\textcolor{red}{\textbf{[TODO: #1]}}}


\begin{document}

% ============================================================
% COVER PAGE
% ============================================================
\begin{titlepage}
    \centering

    \includegraphics[width=4.5cm]{asu_logo.png}\\[1cm]

    {\Huge \textbf{Content-Aware Pseudo-Color Mapping for
    Scientific/Medical Grayscale Images}}\\[1.5cm]

    {\large \textbf{Advanced Digital Image and Video Processing
    (CSE 681)}}\\[0.5cm]

    {\large \textbf{Milestone 1}}\\[1cm]

    {\small
    \textbf{Project Repo:}
    \href{https://github.com/abdulrahman-ragab/content-aware-pseudocolor}
    {\texttt{GitHub Repository}}
    }\\[1.5cm]

    % Team:
    \begin{tabularx}{0.8\textwidth}{|X|c|}
        \hline
        \textbf{Team Members} & \textbf{ID} \\ \hline
        Marwan Hatem Hosny Mobarak & 2600308 \\ \hline
        Kirolos Thabet Fouad Benjamin & 2600316 \\ \hline
        Abdulrahman Mohamed Abdullah Ragab & 2600286 \\ \hline
    \end{tabularx}

    \vspace{1.5cm}

    {\large \textbf{Submitted To:}}\\[0.3cm]

    Dr. Hazem Abbass\\
    Dr. Mahmoud Khalil

    \vfill

\end{titlepage}

\newpage

% ============================================================
% TITLE & MAIN DOCUMENT
% ============================================================

\title{
    \LARGE Project Proposal \& Literature Review \\[6pt]
    \large Content-Aware Pseudo-Color Mapping for
    Scientific/Medical Grayscale Images
}

\author{
    Marwan Hatem Hosny Mobarak
    \and
    Kirolos Thabet Fouad Benjamin
    \and
    Abdulrahman Mohamed Abdullah Ragab
}

\date{
    Milestone 1 --- 2 November 2026 \\[2pt]
    CSE 681 --- Digital Image Processing
}

\maketitle

% ============================================================
% ABSTRACT
% ============================================================

\begin{abstract}

3--5 sentences: what problem are you tackling, why does it matter,
what is your planned approach, and what public dataset/baseline
will you use. Write this last.

\end{abstract}

\noindent
\todo{First Version abstract: revise after completing the
implementation and evaluation experiments to reflect the actual
methods and findings.}

% ============================================================
% 1. PROBLEM STATEMENT & MOTIVATION
% ============================================================

\section{Problem Statement \& Motivation}
\label{sec:motivation}

What image-processing problem are you solving? Why is it interesting
or useful? How does it connect to material from the course
(name the week/topic)? 1 paragraph.

% ============================================================
% 2. RELATED WORK
% ============================================================

\section{Related Work}
\label{sec:related}

Summarize at least 4 references, at least 2 of which should be
from 2023--2026. Use the table below as a starting organizer,
then write 1--2 paragraphs of prose connecting them to your plan.

\begin{table}[htbp]
\centering
\caption{Key references.}
\label{tab:refs}

\begin{tabularx}{\textwidth}{
    @{}
    p{0.30\linewidth}
    p{0.12\linewidth}
    X
    @{}
}
\toprule
\textbf{Reference} &
\textbf{Year} &
\textbf{Relevance to this project} \\
\midrule

\citep{example2024} & 2024 & \todo{one sentence} \\

\citep{example2018} & 2018 & \todo{one sentence} \\

\bottomrule
\end{tabularx}
\end{table}

% ============================================================
% 3. DATASET & BASELINE PLAN
% ============================================================

\section{Dataset \& Baseline Plan}
\label{sec:data}

Name the public dataset(s), give exact size/subset you will use,
and the license. State exactly which baseline (classical or
prior learned method) you will reproduce and compare against.

\begin{table}[htbp]
\centering
\caption{Dataset summary.}
\label{tab:data}

\begin{tabularx}{\textwidth}{
    @{}
    p{0.22\linewidth}
    X
    p{0.32\linewidth}
    @{}
}
\toprule
\textbf{Dataset} &
\textbf{Subset Used} &
\textbf{Link} \\
\midrule

\todo{name} &
\todo{e.g., 800 train / 100 val crops} &
\todo{URL} \\

\bottomrule
\end{tabularx}
\end{table}

% ============================================================
% 4. PROPOSED METHOD
% ============================================================

\section{Proposed Method}
\label{sec:method}

Sketch your planned approach: architecture diagram or pseudocode,
what is novel or extended relative to the baseline, and why you
expect it to work.

% ============================================================
% 5. FEASIBILITY & COMPUTE PLAN
% ============================================================

\section{Feasibility \& Compute Plan}
\label{sec:feasibility}

State your compute tier (CPU-OK / GPU-light), expected training
time per experiment, and confirm this fits a laptop-class machine
within the project window.

% ------------------------------------------------------------
% 5.1 Minimum Viable Deliverable
% ------------------------------------------------------------

\subsection{Minimum Viable Deliverable (MVP)}

The smallest version of this project that still produces a valid,
reportable comparison against the baseline.

% ------------------------------------------------------------
% 5.2 Stretch Goal
% ------------------------------------------------------------

\subsection{Stretch Goal}

\todo{What you will attempt if the MVP is on track by Week 5--6.}

% ============================================================
% 6. RISK ASSESSMENT
% ============================================================

\section{Risk Assessment}
\label{sec:risk}

\begin{itemize}[leftmargin=*]

    \item \todo{Risk 1 and its mitigation / fallback.}

    \item \todo{Risk 2 and its mitigation / fallback.}

\end{itemize}

% ============================================================
% 7. WORK BREAKDOWN
% ============================================================

\section{Work Breakdown}
\label{sec:workbreakdown}

\begin{table}[htbp]
\centering
\caption{Milestone task allocation.}
\label{tab:tasks}

\begin{tabularx}{\textwidth}{
    @{}
    p{0.12\linewidth}
    X
    X
    X
    @{}
}
\toprule

\textbf{Milestone} &
\textbf{Student A} &
\textbf{Student B} &
\textbf{Student C} \\

\midrule

M1 (Wk 3) &
\todo{task} &
\todo{task} &
\todo{task} \\

M2 (Wk 6) &
\todo{task} &
\todo{task} &
\todo{task} \\

M3 (Wk 9) &
\todo{task} &
\todo{task} &
\todo{task} \\

\bottomrule
\end{tabularx}
\end{table}

% ============================================================
% REFERENCES
% ============================================================

\bibliographystyle{unsrtnat}
\bibliography{references}

\end{document}
